\documentclass[11pt]{article}

\usepackage[a4paper,margin=1in]{geometry}
\usepackage{amsmath,amssymb}
\usepackage{enumitem}
\usepackage{hyperref}

\title{Computational Test for Super-Rough Random Dimer Heights}
\author{Alberto Rescigno}
\date{June 2026}

\begin{document}
\maketitle

\begin{abstract}
This note documents the current computational probability project.  The code
uses loop-erased random-walk winding as a tree-side proxy for a dimer height
function under Temperley's bijection.  The numerical target is to distinguish
ordinary rough behaviour,
\[
  \operatorname{Var}(h(0)) \approx C\log L,
\]
from the physics prediction of a super-rough random dimer phase,
\[
  \operatorname{Var}(h(0)) \approx C(\log L)^2 .
\]
The repository now implements quenched random edge-weight environments, all-edge
randomisation, bootstrap scaling diagnostics, and error-bar plots for the main
finite-size probe.
\end{abstract}

\section{Mathematical Setup}

For each box size \(L\), simulations take place in
\[
  B_L=[-L,L]^2\cap \mathbb{Z}^2 .
\]
Every nearest-neighbour edge exposed by the walk receives a positive random
weight.  This is a quenched environment: the edge weight is sampled once and
then kept fixed for the rest of that trial.  The random walk starts at the
origin and chooses neighbouring vertices with probability proportional to the
four incident edge weights.  It stops at the first boundary hit of \(B_L\).

The raw walk is chronologically loop-erased.  The resulting path is interpreted
as the spanning-tree branch corresponding to the origin.  The recorded
observable is the quarter-turn winding
\[
  W_L=\#\{\text{left turns}\}-\#\{\text{right turns}\}.
\]
In the present code,
\[
  h_{\mathrm{proxy}}(0)=W_L .
\]
A final dimer-height convention may multiply this by a fixed
orientation-dependent constant.  Such a normalization changes fitted constants
but not the comparison between \(\log L\) and \((\log L)^2\) growth.

\section{Connection to Dimers}

Temperley's bijection relates dimer coverings on appropriate planar lattice
graphs to spanning trees.  In the Temperleyan setting, dimer height information
can be read through the topology and winding of the associated tree branches.
This is the computational bridge used here: instead of sampling full dimer
coverings directly, the code samples the random-walk side and measures the
winding of the loop-erased branch.

For homogeneous critical dimer models, height fluctuations are expected to be
Gaussian free field type, so the variance at a point grows logarithmically in
the box size.  The research question is whether quenched random weights can
push the same height observable into a super-rough regime with squared-log
growth.

\section{Implemented Random Environments}

The old two-random-weight model
\[
  (w_N,w_E,w_S,w_W)=(1,1,u,v)
\]
has been removed.  It fixed two directions and randomised only two, which was
not the clean random dimer-weight experiment now needed.

The current code randomises all nearest-neighbour edge weights.  Implemented
families include Gamma, Exponential, Lognormal, Pareto, Uniform, Beta, Weibull,
inverse-Gamma, Bernoulli two-point, and triangular distributions.  The simple
symmetric random walk remains as the baseline calibration model.

\section{Statistical Outputs}

Each run writes:
\begin{itemize}[leftmargin=1.5em]
  \item \texttt{raw\_trials.csv}, with one row per walk;
  \item \texttt{summary.csv}, with sample means, variances, and variance
        standard errors for each model and \(L\);
  \item \texttt{fits.csv}, comparing \(C\log L\) and \(C(\log L)^2\), including
        SSE, \(R^2\), AIC, BIC, and bootstrap confidence intervals;
  \item \texttt{run\_report.md}, which records the symmetric baseline, any
        possible super-rough candidates, and the interpretation caveat;
  \item PNG figures with error bars when the plotting libraries are installed.
\end{itemize}

The finite-size model comparison is descriptive.  A BIC preference for
\(C(\log L)^2\) is treated as a numerical signal for further investigation, not
as proof of an asymptotic exponent.

\section{Annealed and Quenched Convention}

The present experiments use one random environment and one walk per trial.
Therefore the reported variance is annealed over both sources of randomness:
the environment and the walk conditional on the environment.  This is the right
first diagnostic for the physics question, but it should be stated explicitly
in any manuscript.

A strictly quenched study would fix one environment and run several independent
walks in that same environment, then compare the conditional variance with the
environment-to-environment variation.  The current code and report separate the
interpretation clearly so this extension can be added without changing the
meaning of the existing results.

\section{Main Probe}

The large preset is:
\begin{verbatim}
python3 main.py run super_rough_probe --workers 8
\end{verbatim}
The preset uses
\[
  L\in\{16,32,64,128,256,512\},
\]
with 1500 samples per model-size pair and a broad set of random edge-weight
distributions.  The output folder is
\[
  \texttt{results/super\_rough\_probe/}.
\]

For the completed publication check on this machine, the same model/size grid
was run with 200 samples per model-size pair:
\begin{verbatim}
python3 main.py run super_rough_probe \
  --output-name super_rough_probe_200 --samples 200 --workers 8
\end{verbatim}
This produced 18,000 walks in 2394.77 seconds.  The output folder is
\[
  \texttt{results/super\_rough\_probe\_200/}.
\]

The symmetric model is the calibration point: it should favour \(C\log L\)
within finite-size error.  A random edge-weight model is interesting if it is
centred, stable across \(L\), and robustly favours \(C(\log L)^2\) by the fit
diagnostics and bootstrap intervals.

\section{Completed Probe Result}

The 200-sample full-grid probe gave the expected baseline behaviour.  For the
symmetric model,
\[
  C_{\log L}=0.773727,
\]
with bootstrap 95\% interval \([0.704848,0.844425]\), and BIC preferred
\(C\log L\).

No non-baseline random edge-weight model preferred \(C(\log L)^2\) by BIC.
The closest finite-size competitors were:
\begin{itemize}[leftmargin=1.5em]
  \item \texttt{lognormal\_edges:0.5}, with
        \(\Delta\mathrm{BIC}=\mathrm{BIC}_{\log^2}-\mathrm{BIC}_{\log}=0.598\);
  \item \texttt{gamma\_edges:1}, with
        \(\Delta\mathrm{BIC}=3.015\);
  \item \texttt{bernoulli\_edges:0.8}, with
        \(\Delta\mathrm{BIC}=4.421\).
\end{itemize}
All of these still favoured \(C\log L\), since positive \(\Delta\mathrm{BIC}\)
means the logarithmic fit had the smaller BIC.

At \(L=512\), the largest observed variance was the symmetric baseline,
\[
  \operatorname{Var}(W_{512})=5.9919 \quad \text{with SE }0.6007.
\]
The largest non-baseline value was \texttt{lognormal\_edges:0.5},
\[
  \operatorname{Var}(W_{512})=5.3216 \quad \text{with SE }0.5335.
\]
Thus, in this completed finite-size run, random edge disorder did not produce
numerical evidence for squared-log height/winding variance.  This is a useful
negative result, not a proof that super-roughness cannot occur in another
normalization, geometry, disorder class, or larger asymptotic range.

\section{Related Work}

Kenyon's work on dimer height fluctuations is the standard reference point for
Gaussian free field behaviour in homogeneous dimer models.  Temperleyan
connections between dimers, spanning trees, and height observables underlie the
winding proxy used here.  Berestycki, Laslier, and Ray give a modern
winding/GFF perspective for dimer height fluctuations through UST branches.

Recent random-environment dimer work motivates the numerical experiments:
\begin{itemize}[leftmargin=1.5em]
  \item Moulard and Toninelli study square-grid dimers with layered quenched
        disorder and show that disorder can strongly change correlations.
  \item Bufetov, Petrov, and Zografos study Aztec-diamond tilings with random
        one-periodic weights, separating GFF-type fluctuations from larger
        environment-driven fluctuations.
  \item Duits and Van Peski introduce Gamma-disordered Aztec-diamond weights and
        connect the model to super-rough random dimer predictions.
  \item Zografos distinguishes quenched and annealed CLTs for one-periodic
        Aztec diamonds in random environment, which is directly relevant to
        how the present simulations should be interpreted.
\end{itemize}

\section{Publication Checklist Status}

\begin{enumerate}[leftmargin=1.6em]
  \item Height convention: implemented as quarter-turn LERW winding, with the
        fixed dimer-height normalization left explicit.
  \item Symmetric calibration: automatically reported in each run's
        \texttt{run\_report.md}.
  \item Super-rough probe: configured as \texttt{super\_rough\_probe} and run
        from the command line.
  \item Error bars and bootstrap intervals: written to \path{summary.csv},
        \path{fits.csv}, plots, and \path{run_report.md}.
  \item Annealed versus quenched interpretation: stated in the generated report
        and in this note.
  \item Literature comparison: included above for the current manuscript draft.
  \item Publication caveat: all fits are finite-size numerical evidence, not
        proofs of asymptotic scaling.
\end{enumerate}

\begin{thebibliography}{9}

\bibitem{kenyon2000}
R. Kenyon.
\newblock Conformal invariance of domino tiling.
\newblock \emph{Annals of Probability}, 28(2), 2000.

\bibitem{berestycki-laslier-ray}
N. Berestycki, B. Laslier, and G. Ray.
\newblock Dimers and imaginary geometry.
\newblock arXiv:1603.09740.

\bibitem{moulard-toninelli}
Q. Moulard and F. Toninelli.
\newblock Dimers with layered disorder.
\newblock arXiv:2507.11964.

\bibitem{bufetov-petrov-zografos}
A. Bufetov, L. Petrov, and P. Zografos.
\newblock Domino Tilings of the Aztec Diamond in Random Environment and Schur
Generating Functions.
\newblock arXiv:2507.08560.

\bibitem{duits-vanpeski}
M. Duits and R. Van Peski.
\newblock The Gamma-disordered Aztec diamond.
\newblock arXiv:2512.03033.

\bibitem{zografos}
P. Zografos.
\newblock Quenched and Annealed CLTs for the one-periodic Aztec diamond in
random environment.
\newblock arXiv:2510.11846.

\end{thebibliography}

\end{document}
